\section{引言}
一项研究从文献阅读推进到实验和结论，需要在多种信息与方法之间作出连续判断。
找到一篇相关论文之后，研究者还要判断其中的条件能否迁移到当前问题；完成一次计算之后，也要决定结果是否足以支持原先的解释。
这类工作既依赖专业知识，又涉及阅读、实现、比较和修正等连续活动。
大语言模型（large language model，LLM）能够在同一语言接口下调用检索、编程与分析工具，使上述活动有可能被纳入相对统一的自动化过程。
科研智能体研究正是在这一背景下兴起\cite{surveyautonomy,surveyagentic}。

检索增强生成（retrieval-augmented generation，RAG）使模型能够利用训练之外的文献\cite{rag}；
ReAct让模型依据工具返回的观察继续行动\cite{react}；AutoGen则提供了组织多个模型实例及其交互的方式\cite{autogen}。
这些方法解决的是不同层次的问题，却常被同时用于一个科研系统。
检索者、程序编写者、评审者等角色由此进入研究流程，形成了目前常见的科研多智能体架构。
本文中的“代理”与“智能体”同义；角色是否构成有效分工，取决于其信息和操作上的差异。

随着系统能够承担的研究步骤增多，完成流程仍不足以保证结论得到支持。
The AI Scientist能够组织实验并生成论文\cite{scientistjournal}，Virtual Lab将跨学科讨论落实到纳米抗体设计与实验表征\cite{virtual}，
对自动科研的审计却发现，最终报告有时不足以暴露实验过程中的方法学缺陷\cite{pitfalls}。
三类结果并不冲突：它们分别考察自动化覆盖、特定科学任务和过程可靠性。
如果只问系统是否“像科学家”，这些差异很容易被同一个概念掩盖。

已有综述对此给出了不同的整理方式。
Zheng等按工具、分析者和科学家区分自主程度\cite{surveyautonomy}；Wei等将系统能力与学科应用联系起来\cite{surveyagentic}；
郑逸宁等则细分了工具使用的规划、指令生成与结果处理\cite{surveytool}。
本文在这些工作的基础上，围绕三个相互关联的问题展开：系统怎样取得并保存研究依据，多个代理怎样据此协作，现有实验又怎样证明这种协作有用。
相比按框架逐项介绍，这一线索更便于解释方法之间的替代关系，也能把课堂涉及的知识图谱、自主记忆和可信科研放在同一个问题下讨论。

本文采用叙述性综述方法，通过题名与关键词检索、引用追踪以及原文核对选取代表工作，资料范围截至2026年11月9日。
分析包括基础机制、科研系统与评价研究，不覆盖全部多智能体强化学习或机器人群体研究，也不进行跨论文效应合并。
同一项目的不同版本按研究进展引用，例如The AI Scientist的早期报告、第二版与期刊稿不作为相互独立的重复验证。
下文先厘清研究对象，再讨论知识、记忆与协作方法；随后以科研应用和公开评测检验这些方法的解释力，最后给出本文的判断及尚未解决的问题。
